% GENERATED FILE. Edit canonical lessons, then run npm run generate:books.
% canonical-source-hash: fnv1a64:7e540d3c47c0f3a4
% canonical-lessons: BN-C422-daripalla, BN-C422-khati, BN-C422-phordo, BN-C422-grahok, BN-C422-porokh-kora

\chapter{Weighing and Testing}
\label{ch:bn-weighing-and-testing}
\hlchaptermodality{422}

\begin{chapteropening}
\textbf{By the end of this chapter:} \emph{I can watch goods weighed, check they are pure and test them before I buy.}

Complete the last lesson of chapter 422: I can watch goods weighed, check they are pure and test them before I buy.
\end{chapteropening}

\section[\texorpdfstring{\bn{দাঁড়িপাল্লা} (dā̃ṛipāllā)}{দাঁড়িপাল্লা (dā̃ṛipāllā)}]{\bn{দাঁড়িপাল্লা} (dā̃ṛipāllā) — a pair of scales; a balance}

\label{lesson:BN-C422-daripalla}



\begin{quote}
Before the new one: say the Bengali for wholesale, then the Bengali for to stock up.
\end{quote}

\subsection*{You'll want to know: \bn{দাঁড়িপাল্লা}}
\textbf{\bn{দাঁড়িপাল্লা}} — \emph{dā̃ṛipāllā} — \textquotedblleft{}a pair of scales; a balance\textquotedblright{}.

\begin{quote}
\bn{মুদি} \bn{দাঁড়িপাল্লায়} \bn{ডাল} \bn{ওজন} \bn{করলেন।}
\end{quote}

\begin{quote}
\emph{mudi dā̃ṛipāllāy ḍāl ojôn kôrlen.}
\end{quote}

\begin{quote}
The grocer weighed the lentils on the scales.
\end{quote}

The hanging scales of a market stall: a beam, the \bn{দাঁড়ি}, with two pans, the \bn{পাল্লা}.

\begin{cousinweb}
\textbf{\bn{দাঁড়িপাল্লা}} (\emph{dā̃ṛipāllā}) joins \bn{দাঁড়ি}, which CDIAL lists as Bengali \emph{ḍā̃ṛi}, \textquotedblleft{}the beam of a balance\textquotedblright{}, under Sanskrit \emph{daṇḍá}, \textquotedblleft{}a stick\textquotedblright{} (CDIAL 6128), to \bn{পাল্লা}, \textquotedblleft{}a scale-pan\textquotedblright{}, whose own origin is not settled.
\end{cousinweb}

\begin{grammarlens}[title={What you've built}]
One more word for getting around.
\end{grammarlens}

\subsection*{Your turn}
\begin{itemize}
\raggedright
  \item \emph{Say it:} \emph{dā̃ṛipāllā}
  \item \emph{Say it:} \emph{dā̃ṛipāllā}, once more
  \item \emph{Say it:} \emph{mudi dā̃ṛipāllāy ḍāl ojôn kôrlen.} — The grocer weighed the lentils on the scales.
  \item \emph{Recall:} say the Bengali for wholesale, then the Bengali for to stock up, then say \emph{dā̃ṛipāllā} again
\end{itemize}

\begin{tcolorbox}[breakable,colback=teal!4,colframe=teal!35!black,title={Before you move on}]
What does \bn{দাঁড়িপাল্লা} mean? (\textquotedblleft{}A pair of scales; a balance\textquotedblright{}.) Say it once more.
\end{tcolorbox}
\section[\texorpdfstring{\bn{খাঁটি} (khā̃ṭi)}{খাঁটি (khā̃ṭi)}]{\bn{খাঁটি} (khā̃ṭi) — pure; genuine; unmixed}

\label{lesson:BN-C422-khati}



\begin{quote}
Before the new one: say the Bengali for to stock up, then the Bengali for a pair of scales.
\end{quote}

\subsection*{You'll want to know: \bn{খাঁটি}}
\textbf{\bn{খাঁটি}} — \emph{khā̃ṭi} — \textquotedblleft{}pure; genuine; unmixed\textquotedblright{}.

\begin{quote}
\bn{এই} \bn{দোকানের} \bn{মধু} \bn{খাঁটি।}
\end{quote}

\begin{quote}
\emph{ei dokāner môdhu khā̃ṭi.}
\end{quote}

\begin{quote}
The honey in this shop is pure.
\end{quote}

Of food, unmixed with anything cheaper; of a person, honest through and through.

\begin{cousinweb}
\textbf{\bn{খাঁটি}} (\emph{khā̃ṭi}) is the everyday word for pure, unmixed goods; its origin is not settled.
\end{cousinweb}

\begin{grammarlens}[title={What you've built}]
One more word for getting around.
\end{grammarlens}

\subsection*{Your turn}
\begin{itemize}
\raggedright
  \item \emph{Say it:} \emph{khā̃ṭi}
  \item \emph{Say it:} \emph{khā̃ṭi}, once more
  \item \emph{Say it:} \emph{ei dokāner môdhu khā̃ṭi.} — The honey in this shop is pure.
  \item \emph{Recall:} say the Bengali for to stock up, then the Bengali for a pair of scales, then say \emph{khā̃ṭi} again
\end{itemize}

\begin{tcolorbox}[breakable,colback=teal!4,colframe=teal!35!black,title={Before you move on}]
What does \bn{খাঁটি} mean? (\textquotedblleft{}Pure; genuine; unmixed\textquotedblright{}.) Say it once more.
\end{tcolorbox}
\section[\texorpdfstring{\bn{ফর্দ} (phôrdô)}{ফর্দ (phôrdô)}]{\bn{ফর্দ} (phôrdô) — a shopping list; a list}

\label{lesson:BN-C422-phordo}



\begin{quote}
Before the new one: say the Bengali for a pair of scales, then the Bengali for pure (unmixed).
\end{quote}

\subsection*{You'll want to know: \bn{ফর্দ}}
\textbf{\bn{ফর্দ}} — \emph{phôrdô} — \textquotedblleft{}a shopping list; a list\textquotedblright{}.

\begin{quote}
\bn{বাজারে} \bn{যাওয়ার} \bn{আগে} \bn{মীরা} \bn{একটা} \bn{ফর্দ} \bn{লিখল।}
\end{quote}

\begin{quote}
\emph{bājāre jāoyār āge mīrā ekṭā phôrdô likhlô.}
\end{quote}

\begin{quote}
Before going to the market Mira wrote a list.
\end{quote}

\begin{cousinweb}
\textbf{\bn{ফর্দ}} (\emph{phôrdô}) is usually traced to Persian \emph{fard}, \textquotedblleft{}a single sheet, a list\textquotedblright{}, from Arabic.
\end{cousinweb}

\begin{grammarlens}[title={What you've built}]
One more word for getting around.
\end{grammarlens}

\subsection*{Your turn}
\begin{itemize}
\raggedright
  \item \emph{Say it:} \emph{phôrdô}
  \item \emph{Say it:} \emph{phôrdô}, once more
  \item \emph{Say it:} \emph{bājāre jāoyār āge mīrā ekṭā phôrdô likhlô.} — Before going to the market Mira wrote a list.
  \item \emph{Recall:} say the Bengali for a pair of scales, then the Bengali for pure (unmixed), then say \emph{phôrdô} again
\end{itemize}

\begin{tcolorbox}[breakable,colback=teal!4,colframe=teal!35!black,title={Before you move on}]
What does \bn{ফর্দ} mean? (\textquotedblleft{}A shopping list; a list\textquotedblright{}.) Say it once more.
\end{tcolorbox}
\section[\texorpdfstring{\bn{গ্রাহক} (grāhôk)}{গ্রাহক (grāhôk)}]{\bn{গ্রাহক} (grāhôk) — a customer; a subscriber}

\label{lesson:BN-C422-grahok}



\begin{quote}
Before the new one: say the Bengali for pure (unmixed), then the Bengali for a shopping list.
\end{quote}

\subsection*{You'll want to know: \bn{গ্রাহক}}
\textbf{\bn{গ্রাহক}} — \emph{grāhôk} — \textquotedblleft{}a customer; a subscriber\textquotedblright{}.

\begin{quote}
\bn{দোকানে} \bn{নতুন} \bn{গ্রাহকদের} \bn{জন্য} \bn{ছাড়} \bn{আছে।}
\end{quote}

\begin{quote}
\emph{dokāne nôtun grāhôkder jônyô chhāṛ āchhe.}
\end{quote}

\begin{quote}
The shop has a discount for new customers.
\end{quote}

\bn{খদ্দের} is the everyday word; \bn{গ্রাহক} is what banks, shops and the post office say.

\begin{cousinweb}
\textbf{\bn{গ্রাহক}} (\emph{grāhôk}) is Sanskrit \emph{grāhaka}, \textquotedblleft{}one who receives, a purchaser\textquotedblright{} (Monier-Williams), from \emph{grah}, to take.
\end{cousinweb}

\begin{grammarlens}[title={What you've built}]
One more word for getting around.
\end{grammarlens}

\subsection*{Your turn}
\begin{itemize}
\raggedright
  \item \emph{Say it:} \emph{grāhôk}
  \item \emph{Say it:} \emph{grāhôk}, once more
  \item \emph{Say it:} \emph{dokāne nôtun grāhôkder jônyô chhāṛ āchhe.} — The shop has a discount for new customers.
  \item \emph{Recall:} say the Bengali for pure (unmixed), then the Bengali for a shopping list, then say \emph{grāhôk} again
\end{itemize}

\begin{tcolorbox}[breakable,colback=teal!4,colframe=teal!35!black,title={Before you move on}]
What does \bn{গ্রাহক} mean? (\textquotedblleft{}A customer; a subscriber\textquotedblright{}.) Say it once more.
\end{tcolorbox}
\section[\texorpdfstring{\bn{পরখ} \bn{করা} (pôrôkh kôrā)}{পরখ করা (pôrôkh kôrā)}]{\bn{পরখ} \bn{করা} (pôrôkh kôrā) — to test; to try out}

\label{lesson:BN-C422-porokh-kora}



\begin{quote}
Before the new one: say the Bengali for a shopping list, then the Bengali for a customer (formal).
\end{quote}

\subsection*{You'll want to know: \bn{পরখ} \bn{করা}}
\textbf{\bn{পরখ} \bn{করা}} — \emph{pôrôkh kôrā} — \textquotedblleft{}to test; to try out\textquotedblright{}.

\begin{quote}
\bn{কেনার} \bn{আগে} \bn{আমি} \bn{আমটা} \bn{টিপে} \bn{পরখ} \bn{করলাম।}
\end{quote}

\begin{quote}
\emph{kenār āge āmi āmṭā ṭipe pôrôkh kôrlām.}
\end{quote}

\begin{quote}
Before buying, I pressed the mango to test it.
\end{quote}

Past (he/she …): \bn{তিনি} \bn{জিনিসটা} \bn{পরখ} \bn{করলেন}, \emph{tini jinisṭā pôrôkh kôrlen}.

\begin{cousinweb}
\textbf{\bn{পরখ} \bn{করা}} (\emph{pôrôkh kôrā}) joins \bn{পরখ}, which CDIAL lists as Bengali \emph{parakh}, \textquotedblleft{}an examination\textquotedblright{}, under Sanskrit \emph{párīkṣatē}, \textquotedblleft{}observes, examines\textquotedblright{} (CDIAL 7904), to \bn{করা}, to do; \bn{পরীক্ষা}, an exam, is the same word taken whole.
\end{cousinweb}

\begin{grammarlens}[title={What you've built}]
One more word for getting around.
\end{grammarlens}

\subsection*{Your turn}
\begin{itemize}
\raggedright
  \item \emph{Say it:} \emph{pôrôkh kôrā}
  \item \emph{Say it:} \emph{pôrôkh kôrā}, once more
  \item \emph{Say it:} \emph{kenār āge āmi āmṭā ṭipe pôrôkh kôrlām.} — Before buying, I pressed the mango to test it.
  \item \emph{Recall:} say the Bengali for a shopping list, then the Bengali for a customer (formal), then say \emph{pôrôkh kôrā} again
\end{itemize}

\begin{tcolorbox}[breakable,colback=teal!4,colframe=teal!35!black,title={Before you move on}]
What does \bn{পরখ} \bn{করা} mean? (\textquotedblleft{}To test; to try out\textquotedblright{}.) Say it once more.
\end{tcolorbox}
